%======================================================================
%  Companion to Chapter 3 -- Identification of a planar crack by the
%  reciprocity gap
%======================================================================
\chapter{Identification of a planar crack by the reciprocity gap}\label{cc:ch3}

\framepart{Outline}

Chapter~\ref{B-ch:rg} of the notes identifies a planar crack from boundary
data in three steps: the normal, the line (or plane), and the extension. In
two dimensions all three are computed by one Cast3M file on the example of
Section~\ref{B-sec:rg-num}: the rectangle $1.5\times2$ with the crack from
$(0.5,\,1.2)$ to $(0.8,\,1.35)$, of length $0.3354$. The file builds the cracked
mesh, runs two experiments with electrodes, measures the potential and the flux
on the boundary, adds noise, and identifies the crack with the adjoint fields of
the chapter. A Python twin of the file, on the same mesh divisions, gives the
reference values.

You should then be able to build a crack into a Cast3M mesh, compute a
reciprocity gap from boundary nodal values only, say which adjoint fields the
discrete computation represents exactly, and read the effect of the mesh, of
the noise and of the orientation of the loads on each step of the
identification.

\section{Where Cast3M enters Chapter 3}\label{cc:sec3-map}
\index{Cast3M!Chapter 3}%

\begin{gbox}{Chapter 3 of the notes and Cast3M}\label{cc:tab3}
\small
\begin{tabular}{@{}p{42mm}@{\hspace{3mm}}p{98mm}@{}}
Direct problem, Sec.~\ref{B-sec:rg-direct}
 & \textbf{Exercise~\ref{exo:cc-rg-direct}}: the thermal model on a mesh with
   two lips; the two openings against the infinite-plate formula; the body
   without crack, $\RG=0$.\\[1mm]
Reciprocity gap, Sec.~\ref{B-sec:rg-gap}
 & \textbf{Section~\ref{cc:sec3-gap}}: $\RG$ from nodal values, exact for
   adjoint fields of degree at most two on quadratic triangles.\\[1mm]
Steps 1--2, moments, dipole, Sec.~\ref{B-sec:rg-rec}
 & \textbf{Exercises~\ref{exo:cc-rg-steps} and~\ref{exo:cc-rg-mesh}}: normal,
   line, centre, length by the second moment and by the dipole; the cubic
   field on three meshes.\\[1mm]
Stability, Sec.~\ref{B-ssec:rg-stab}; Ex.~\ref{B-exo:rg-codes}
 & \textbf{Exercises~\ref{exo:cc-rg-noise} and~\ref{exo:cc-rg-sine}}: noise on
   the data, dipole against second moment; the sine--sinh series and its
   resolution.\\[1mm]
Several experiments, Sec.~\ref{B-sec:rg-multi}
 & \textbf{Exercise~\ref{exo:cc-rg-rotation}}: the loads turned by $\beta$,
   the blind direction, one experiment against a pair.\\[1mm]
3D cracks, Sec.~\ref{B-sec:rg-rec}
 & To do: a penny-shaped crack in a cube (\texttt{CUB20}), three experiments,
   the dipole estimate $a=(3m_0/8E_n)^{1/3}$.\\[1mm]
Elasticity, Ex.~\ref{B-exo:antiplane}
 & To do: the same file in anti-plane shear is unchanged (the potential is the
   displacement $u_3$); in plane elasticity the gap uses the elastic adjoint
   fields of Andrieux, Ben Abda and Bui.\\
\end{tabular}
\end{gbox}

\section{A crack in a mesh}\label{cc:sec3-mesh}
\index{crack!in a mesh}%
A crack is a line whose two sides do not share their nodes, except at the tips.
In gibiane this is two lines with the same end points: each call of \op{droi}
creates its own interior nodes, and the end points are the same objects.
\begin{castemcom}
The crack $A$--$B$ is drawn twice, \texttt{dplus} and \texttt{dmoins}. The two
lines \texttt{d1a} and \texttt{db3} prolong it to the corners $p_1$ and
$p_3$, and cut the rectangle into two parts, meshed separately by \op{surf}.
They share the nodes of \texttt{d1a}, \texttt{db3} and of the tips, and not
those of the crack: the union is the cracked body.

Meshing the upper part on \texttt{dmoins} instead of \texttt{dplus} gives the
same body without crack, which the file uses as a check.
\end{castemcom}
\begin{castemlines}
\begin{verbatim}
pa = xpa ypa;  pb = xpb ypb;
dplus  = droite (nelem/2) pa pb;
dmoins = droite (nelem/2) pa pb;
d1a = droite (nelem/2) p1 pa;
db3 = droite (nelem/2) pb p3;
smoins = surf (d12 et d23 et db3
               et dmoins et d1a);
splus  = surf (d1a et dplus et db3
               et d34 et d41);
dom  = smoins et splus;
\end{verbatim}
\end{castemlines}

\section{The reciprocity gap from nodal values}\label{cc:sec3-gap}
\index{reciprocity gap!discrete}%
With the conductivity matrix $\mat{K}$ of the cracked mesh and the nodal
potential $\mat{u}$ of an experiment, the nodal fluxes are $\mat{q}=\mat{K}\mat{u}$:
zero at the free nodes, the reactions of the electrodes at the others. An
adjoint field $v$ is interpolated at the nodes, with the same value on both
lips. The file computes, on the nodes of the outer boundary only,
\begin{equation}\label{cc:eq-rgh}
  \RG_h(v)=\sum_{i\in\partial\Omega}\bigl(u_i\,(\mat{K}\mat{v})_i-v_i\,q_i\bigr),
\end{equation}
where $(\mat{K}\mat{v})_i$ is the nodal flux of $v$. The procedure
\texttt{rgap} of the file returns $-\RG_h(v)$; the sign cancels in all the
ratios of the identification.

\begin{gbox}{The discrete reciprocity gap is exact for adjoint fields of degree two}\label{cc:box-rgh}
On a mesh of quadratic triangles, $\displaystyle \RG_h(v)=\int_\Gamma\jump{u_h}\,\partial_nv\,ds$
exactly, for every harmonic polynomial $v$ of degree at most two. The normal,
the line, $m_0$, the centre and the dipole length are therefore exact functions
of the discrete opening, on any mesh. The cubic field of the second moment and
the sine--sinh fields are not represented, and their gaps carry a
discretization error.
\end{gbox}

\paragraph{Why it works.}
Split the boundary sums of~\eqref{cc:eq-rgh} into sums over all nodes. Since
$\mat{q}=\mat{K}\mat{u}$ vanishes away from the outer boundary,
$\sum_{\partial\Omega}v_iq_i=\mat{v}\T\mat{K}\mat{u}$. For $v$ harmonic and of
degree at most two, $v$ is its own interpolant, and
$\displaystyle (\mat{K}\mat{v})_i=\int\nabla v\cdot\nabla N_i$ is the flux of
$v$ through the boundary of the support of $N_i$: zero at the interior nodes,
$\displaystyle \int_{\partial\Omega}\partial_\nu v\,N_i$ on the outer boundary, and
$\displaystyle \mp\int_{\Gamma^\pm}\partial_nv\,N_i$ at the nodes of the lips. Hence
$\sum_{\partial\Omega}u_i(\mat{K}\mat{v})_i=\mat{u}\T\mat{K}\mat{v}
+\int_{\Gamma}\jump{u_h}\,\partial_nv$, and the symmetry of $\mat{K}$ cancels
$\mat{u}\T\mat{K}\mat{v}-\mat{v}\T\mat{K}\mat{u}$. For a cubic $v$ the
interpolant differs from $v$, the interior fluxes $(\mat{K}\mat{v})_i$ no longer
vanish, and the difference is an error of order $h^2$.

\begin{castemcom}
The procedure works on boundary data only: \texttt{ub} and \texttt{qb} live on
the outer boundary, and every field is reduced to that support before the
scalar products \op{xty}. \texttt{kk * vt} is the nodal flux of the adjoint
field.
\end{castemcom}
\begin{castemlines}
\begin{verbatim}
debproc rgap ub*chpoint qb*chpoint
             v*chpoint kk*rigidite;
  mb = extr ub 'MAIL';
  vt = exco v 'SCAL' 'T';
  qv = redu mb (kk * vt);
  vb = redu mb vt;
  r = (xty qb vb (mots 'Q') (mots 'T'))
    - (xty ub qv (mots 'T') (mots 'Q'));
finproc r;
\end{verbatim}
\end{castemlines}

\framepart{Summary}

\begin{gbox*}
\begin{itemize}
\item \textbf{A crack is two lines with the same ends}: separate interior
  nodes, shared tips; the body without crack is the same mesh with one lip.
\item \textbf{The gap needs only boundary nodal values}: potentials, the
  nodal fluxes $\mat{K}\mat{u}$, and the nodal fluxes $\mat{K}\mat{v}$ of the
  adjoint field.
\item \textbf{Degree two is exact}: on quadratic triangles the normal, the
  line, $m_0$, the centre and the dipole length are exact functions of the
  discrete opening; the second moment and the sine--sinh series carry a
  discretization error of order $h^2$ that is visible without noise.
\item \textbf{The dipole estimate is the robust length}: it uses $m_0$ only,
  and survives $1\%$ noise; the second moment does not.
\item \textbf{Two loads remove the blind direction}: a single uniform field
  parallel to the crack sees nothing; the optimal combination of two
  orthogonal loads sees the crack at every angle.
\end{itemize}
\end{gbox*}

\framepart{Exercises}

The exercises use the file \texttt{crack\_rg\_extension.dgibi} and its Python
twin \texttt{rg\_reference.py}, which builds the same cracked rectangle with
the same divisions (30 per side, 15 on the crack and on the two cut lines),
quadratic triangles, and computes the same quantities. The numbers below come
from the twin; the Cast3M mesh of \op{surf} is different, so expect agreement to
the second or third digit, and to all digits for the normal and the line.

\framesub{Short questions}

\exercise{The sign of the gap}\label{exo:cc-rg-q-sign}
\index{reciprocity gap!sign}%
The procedure \texttt{rgap} computes $\sum(q\,v-u\,\mat{K}v)$, that is
$-\RG(v)$. Which results of the identification change sign, and which do not?
\begin{solution}
$m_0$, $m_1$ and the raw normal $(\RG(x),\RG(y))$ change sign. The normal is
defined up to its sign anyway, and $c=\RG(v_2)/(2\RG(\eta))$,
$s_0=m_1/m_0$, $\mu_2=\RG(v_3)/m_0$ and $a=\sqrt{\abs{m_0}/(\pi\abs{E_n})}$ are
ratios or absolute values: they do not change. The check without crack, where
every gap vanishes, does not see the sign either; the flux balance
$\RG(1)=-\sum q_i=0$ checks the fluxes.
\end{solution}

\exercise{The blind direction of the rectangle}\label{exo:cc-rg-q-blind}
\index{blind experiment}%
The potential $u=\vd\cdot\vx$ is imposed on the whole boundary, with
$\vd=(\cos\beta,\sin\beta)$. For which $\beta$ is the experiment blind, and for
which is the signal largest?
\begin{solution}
Blind when $\vd$ is parallel to the crack, $\vd\cdot\vn=0$:
$\tan\beta=0.15/0.3$, $\beta=26.57^\circ$ (and $206.57^\circ$). The signal
$m_0\propto\vd\cdot\vn$ is largest for $\vd=\pm\vn$, $\beta=116.57^\circ$. The
twin gives $\abs{m_0}=0.0374$ at $0^\circ$, $0.0023$ at $25^\circ$ and the
maximum $0.0836$ near $115^\circ$.
\end{solution}

\exercise{Exact or not?}\label{exo:cc-rg-q-exact}
\index{adjoint field!polynomial}%
On quadratic triangles, which of the adjoint fields $x$, $y$,
$(\vn\cdot\vx)^2-(\vt\cdot\vx)^2$, $\eta$, $s\eta$, $s^2\eta-\tfrac13\eta^3$ and
$\sin(\lambda s)\sinh(\lambda\eta)$ give an exact discrete gap?
\begin{solution}
The first five, of degree at most two (Section~\ref{cc:sec3-gap}). The cubic
field of the second moment and the sine--sinh fields are not in the
finite element space: their gaps carry an error of order $h^2$, growing with
$\lambda$ for the sine--sinh fields.
\end{solution}

\framesub{Problems}

\exercise{The cracked rectangle}\label{exo:cc-rg-direct}
\index{direct problem!cracked body}%
Mesh the rectangle $1.5\times2$ with the crack from $(0.5,1.2)$ to
$(0.8,1.35)$. Experiment E1 imposes $u=0$ at the bottom and $u=2$ at the top
(background field $\vE=(0,1)$), experiment E2 imposes $u=0$ on the left and
$u=1.5$ on the right ($\vE=(1,0)$); the other sides are insulated. (a)~Compute
the openings of E1 and E2 along the crack and compare them with the
infinite-plate opening $2E_n\sqrt{a^2-s^2}$ of Section~\ref{B-sec:rg-num}.
(b)~Check that the gaps vanish for the body without crack.
\begin{solution}
(a)~Both openings are close to half-ellipses, with $m_0=0.0754$ for E1 and
$-0.0376$ for E2, in the ratio $\vE_1\cdot\vn:\vE_2\cdot\vn=0.894:(-0.447)$. On
the middle $80\%$ of the crack the opening of E1 is $0.94$ to $0.97$ times the
infinite-plate value $2E_n\sqrt{a^2-s^2}$ ($0.974$ at the centre); with 120
divisions per side, $0.98$ to $0.99$: the rest is the finite body. At the nodes
next to the tips the ratio falls to $0.8$, because the square-root singularity
is not resolved. The second-moment length of both true openings is $0.328$,
below the true length $0.335$ for the same reason.
(b)~Without crack $\RG(x)$ and $\RG(y)$ vanish to rounding; with the crack the
flux balance gives $\RG(1)=8\times10^{-14}$.
\end{solution}

\exercise{The three steps, from one and from two experiments}\label{exo:cc-rg-steps}
\index{crack!planar, reconstruction}%
With exact data, identify the crack from E1 alone and from the optimal
combination of E1 and E2 (Section~\ref{B-sec:rg-multi}): normal, line, centre,
length by the second moment and by the dipole, and the error on the tips.
\begin{solution}
Both give the normal $\pm(0.4472,-0.8944)$ and the line $c=\pm0.2907$ exactly,
as the box of Section~\ref{cc:sec3-gap} says. E1: $\abs{m_0}=0.0756$,
second-moment length $0.329$, dipole length $0.328$, tips within $0.004$. The
combination has the weights $(0.895,-0.447)$, i.e.\ the background field
normal to the crack ($E_n=1$), and $\abs{m_0}=0.0845$; its dipole length is
again $0.328$, but its second-moment length drops to $0.302$: the cubic field is
not exact, and its error is large for E2 on this mesh
(Exercise~\ref{exo:cc-rg-mesh}). The book's values on a coarser Cast3M mesh,
$m_0=0.0767$ and $0.330$, differ by the mesh.
\untested{\texttt{cast3m/ch3/crack\_rg\_extension.dgibi} (version 3)}
\lstinputlisting[style=gibstyle,firstline=137,lastline=277]{cast3m/ch3/crack_rg_extension.dgibi}
\end{solution}

\exercise{The second moment on three meshes}\label{exo:cc-rg-mesh}
\index{moments of the opening!discretization error}%
Compute the second-moment length from E2 alone, with the exact line and centre,
on meshes with 30, 60 and 120 divisions per side. Explain why it converges,
and compare with the dipole length.
\begin{solution}
$0.147$, $0.300$ and $0.326$ (E1: $0.329$, $0.333$, $0.334$): the error is
divided by about 5 and then 4, an $h^2$ convergence towards the continuous
length $0.335$. The cubic field $s^2\eta-\tfrac13\eta^3$ is not in the space of
quadratic elements, so its gap carries the error of
Section~\ref{cc:sec3-gap}; E2 has half the signal of E1 and suffers more. The
dipole length, $0.328$, $0.331$, $0.333$, is the same for E1 and E2 to
$0.001$ on every mesh: it uses $m_0=\RG(\eta)$, exact on every mesh, and changes
only with the discrete opening itself. The twin reads the number of divisions from the
variable \texttt{RG\_NELEM}.
\lstinputlisting[firstline=8,lastline=97]{cast3m/ch3/rg_reference.py}
\end{solution}

\exercise{Noise: the dipole against the second moment}\label{exo:cc-rg-noise}
\index{dipole}\index{stability!of the extension step}%
Add a relative Gaussian noise $\epsilon$ to the measured potentials and fluxes
(not to the imposed ones), 200 times, for $\epsilon=10^{-3}$ and $10^{-2}$, and
compare the lengths given by the dipole and by the second moment
(Exercise~\ref{B-exo:rg-codes}).
\begin{solution}
At $\epsilon=10^{-3}$ (the default of the file), E1 gives the dipole length
$0.328\pm0.001$ and the second-moment length $0.325\pm0.073$, with 3 failures
($\mu_2\le0$) in 200. At $\epsilon=10^{-2}$: dipole $0.328\pm0.010$, second
moment $0.61\pm0.22$ with 77 failures. The combination of E1 and E2 reduces the
dipole scatter to $\pm0.007$. The dipole uses only $m_0$, measured with a linear
field; the second moment needs $\mu_2m_0$, of order $a^4$, through a field that
contains $\eta^3$. The angle error of the normal is $0.25^\circ$ and
$2.6^\circ$ at the two noise levels.
\end{solution}

\exercise{The sine--sinh series}\label{exo:cc-rg-sine}
\index{sine--sinh adjoint fields}\index{resolution}%
Compute the coefficients $\RG(v_n)$, $n=1,\dots,7$, of the series of
Section~\ref{B-sec:rg-num} for E1 with exact data, along the identified line,
with $L=\sqrt2\times2=2.83$. Plot the partial sums. Which terms can be used, and
what does the series tell about the crack?
\begin{solution}
$\RG(v_n)=0.054$, $0.075$, $0.054$, $0.017$, $-0.021$, then $-0.38$ and
$-2.30$: the last two are discretization errors of the sine--sinh fields, which
reach $\sinh(\lambda_7H)\approx10^5$ on the boundary and are not represented by
elements of size $0.05$. Up to $n=5$ the partial sums show one bump near the
crack, peaking at $g=0.55$ to $0.65$ (the crack covers $0.56\le g\le0.89$,
centre $0.73$), of area $0.09$ (against $m_0=0.076$) and maximum $0.11$ to $0.12$
(against $0.29$ for the true opening). With $H=1.52$ and $\delta=10^{-3}$ the estimate~\eqref{B-eq:rg-res}
gives about four usable terms and a resolution of $0.7$, twice the length of the
crack: the series locates the crack and measures its opening area, not its
shape. Version 2 of the file plotted $n=7$; version 3 stops at $n=5$.
\end{solution}

\exercise{Turning the loads}\label{exo:cc-rg-rotation}
\index{several experiments}\index{virtual experiment}%
Impose $u=\vd\cdot\vx$ on the whole boundary, $\vd=(\cos\beta,\sin\beta)$,
measure only the flux, with $\epsilon=10^{-3}$, and identify the crack for
$\beta=0^\circ,5^\circ,\dots,180^\circ$ from this experiment alone and from the
pair $(\beta,\beta+90^\circ)$ combined.
\begin{solution}
With exact data both identify the normal exactly at every angle; only the
signal $\abs{m_0}$ changes, from $0.0023$ at $25^\circ$ to $0.0836$. With noise,
the single experiment loses the crack near the blind direction: angle errors
of $1$ to $2^\circ$ between $10^\circ$ and $35^\circ$, and a tip error of $0.14$ at
$25^\circ$. The pair keeps angle errors below $0.3^\circ$ and tip errors below
$0.008$ at every angle: its optimal combination is always the load normal to the
crack (Section~\ref{B-sec:rg-multi}). By linearity the file computes the two
fields $u=x$ and $u=y$ once and combines them for every $\beta$.
\lstinputlisting[style=gibstyle,firstline=473,lastline=536]{cast3m/ch3/crack_rg_extension.dgibi}
\end{solution}

\begin{chapterbib}{9}
\bibitem{cc-ab96} S.~Andrieux, A.~Ben Abda, Identification of planar cracks by
  complete overdetermined data: inversion formulae, \emph{Inverse Problems} 12
  (1996) 553--563.
\bibitem{cc-ac09} A.~Constantinescu, \texttt{crack\_reciprocity\_plan.dgibi},
  Cast3M procedure, CNRS -- \'Ecole Polytechnique, 2009.
\end{chapterbib}
